

\TNchapter[Agentic systems represent a shift from static automation to adaptive, goal-directed intelligence embedded across the enterprise.]{Safety and Ethical Evaluation}
\section{Harmful Output Detection }
\subsection{Introduction}
Harmful output detection assesses whether agents produce content or take actions that damage users/systems or violate ethical boundaries. This shifts focus from text harmlessness to action outcomes.

\subsection{Evaluation Dimensions}
\begin{enumerate}
    \item \textbf{Content Safety}: Detecting toxic or inappropriate language.
    \item \textbf{Behavioral Safety}: Identifying risky actions (e.g., file deletion) and tool misuse.
    \item \textbf{Interaction Safety}: Vulnerability to prompt injection or goal hijacking.
    \item \textbf{System Safety}: Compliance with privacy and security constraints.
\end{enumerate}

\subsection{Tools and Frameworks}
\begin{itemize}
    \item \textbf{Azure AI Foundry}: \texttt{ContentSafetyEvaluator} (harm categories), \texttt{IndirectAttackEvaluator} (prompt injection), \texttt{CodeVulnerabilityEvaluator}.
    \item \textbf{Benchmarks}: \textbf{Agent-SafetyBench} (interaction risks), \textbf{AgentHarm} (malicious compliance/jailbreaking), \textbf{R-Judge} (safety awareness).
\end{TNtwocol}
\begin{TNtwocol}
\end{itemize}
\section{Bias and Fairness Assessment}
\subsection{Introduction}
Bias assessment ensures agents behave equitably across diverse populations, addressing representational, performance, and allocative disparities.

\subsection{Sources and Manifestations}
Bias arises from training data, RLHF, or tool outputs. It manifests as:
\begin{itemize}
    \item \textbf{Representation Bias}: Stereotypical or imbalanced content generation.
    \item \textbf{Performance Disparities}: Differential accuracy across demographics (e.g., dialects).
    \item \textbf{Allocative Bias}: Unequal distribution of resources (e.g., loan approvals).
\end{itemize}

\subsection{Evaluation Methodology}
\begin{enumerate}
    \item \textbf{Demographic Variation Testing}: Varying names/pronouns while holding tasks constant.
    \item \textbf{Edge Case Analysis}: Testing non-standard dialects or intersectional identities.
    \item \textbf{Adversarial Testing}: Probing for latent prejudices.
\end{enumerate}
\end{TNtwocol}
\begin{TNtwocol}
\section{Demographic Parity and Equal Opportunity Metrics}
\subsection{Foundational Concepts}
\begin{enumerate}
    \item \textbf{Demographic Parity}: Requires outcome independence from group membership (e.g., equal acceptance rates). Useful when ground truth is unavailable but ignores qualification differences.
    \item \textbf{Equal Opportunity}: Requires equal True Positive Rates (sensitivity) across groups. Focuses on qualified individuals having equal chances.
    \item \textbf{Equalized Odds}: Requires parity in both True Positive and False Positive rates.
\end{enumerate}

\end{TNtwocol}
\begin{TNtwocol}
\subsection{Measurement}
Key metrics include Selection Rate differences (Demographic Parity Difference) and True Positive Rate differences (Equal Opportunity Difference). Statistical significance testing is crucial to rule out random variation.
\section{Toxicity and Offensive Content Screening }
\subsection{Introduction}
Toxicity screening prevents agents from generating hate speech, harassment, or unsafe content. Unlike filtration, it must proactively detect harmful outputs before delivery.

\subsection{Toxicity Types}
Includes Hate Speech, Harassment, Profanity, Threats, Identity Attacks, and Sexually Explicit Content. Toxicity can be explicit or implicit (coded language).

\subsection{Evaluation}
\begin{itemize}
    \item \textbf{Explicit Detection}: Using benchmarks like \textbf{RealToxicityPrompts} and \textbf{CivilComments}.
    \item \textbf{Implicit/Contextual}: Identifying context-dependent microaggressions.
    \item \textbf{False Positive Analysis}: Ensuring legitimate discussions (e.g., about discrimination) aren't flagged.
\end{itemize}
\end{TNtwocol}
\begin{TNtwocol}
\section{Privacy Leakage Detection }
\subsection{Introduction}
Privacy leakage involves inadvertent disclosure of sensitive data through memorization, inference, or tool outputs.

\subsection{Leakage Categories}
\begin{itemize}
    \item \textbf{Training Data Memorization}: Regurgitating PII encountered during training.
    \item \textbf{PII Disclosure}: Direct (SSN, email) or indirect identifiers.
    \item \textbf{Cross-User Contamination}: Leaking one user's context to another in multi-tenant systems.
    \item \textbf{Attribute Inference}: Deducing protected attributes (health, political) from inputs.
\end{itemize}

\end{TNtwocol}
\begin{TNtwocol}
\subsection{Detection Tools}
Automated PII scanners (Microsoft Presidio, LLM Guard) measure PII Exposure Rate and Precision/Recall of leakage detection. Attempts include Passive Inference and Active Extraction (jailbreaking).
\section{Ethical Boundary Compliance }
\subsection{Introduction}
Ethical compliance ensures agents respect values and avoid prohibited behaviors (e.g., manipulation, deception) even when technically feasible.

\subsection{Foundational Principles}
\begin{itemize}
    \item \textbf{Autonomy}: Respecting user self-determination and consent.
    \item \textbf{Non-Maleficence}: Avoiding harm (illegal acts, dangerous instructions).
    \item \textbf{Justice}: Equitable treatment across groups.
    \item \textbf{Beneficence}: Promoting well-being without paternalism.
\end{itemize}

\subsection{Evaluation Framework}
Assessment involves \textbf{Moral Dilemma Scenarios} (testing reasoning in trade-offs) and \textbf{Boundary Testing} (measuring refusal rates for unethical requests). Agents are judged on both decision outcomes and the quality of ethical reasoning.
\end{TNtwocol}

